% BHU PYQ -- Manually transcribed & error-corrected
% Paper: MTB-509 : Special Theory of Relativity-I
% Exam: B.A./B.Sc. (Hons) Semester V Examination, 2022-23

\documentclass[11pt,a4paper]{article}
\usepackage[a4paper,top=2.5cm,bottom=2.5cm,left=2.8cm,right=2.8cm,headheight=14pt]{geometry}
\usepackage{amsmath,amssymb}
\usepackage[T1]{fontenc}\usepackage[utf8]{inputenc}
\usepackage{lmodern,microtype,fancyhdr,enumitem,hyperref,lastpage,parskip}

\hypersetup{colorlinks=true,urlcolor=black,linkcolor=black,
  pdftitle={MTB-509: Special Theory of Relativity-I -- BHU B.A./B.Sc. Sem V 2022-23}}
\pagestyle{fancy}\fancyhf{}
\renewcommand{\headrulewidth}{0.4pt}\renewcommand{\footrulewidth}{0.4pt}
\fancyhead[L]{\small\textit{Banaras Hindu University}}
\fancyhead[R]{\small\textit{MTB-509: Relativity-I}}
\fancyfoot[C]{\small Page \thepage\ of \pageref{LastPage}}

\newlist{parts}{enumerate}{2}
\setlist[parts,1]{label=(\alph*),leftmargin=*,itemsep=5pt,topsep=3pt}
\setlist[parts,2]{label=(\roman*),leftmargin=*,itemsep=3pt,topsep=2pt}
\newcommand{\pts}[1]{\hfill\textit{[#1]}}

\begin{document}
\begin{center}
  {\large\bfseries BANARAS HINDU UNIVERSITY}\\[2pt]
  {\normalsize B.A./B.Sc. (Hons) --- Semester V Examination, 2022-23}\\[6pt]
  \rule{0.8\linewidth}{0.5pt}\\[6pt]
  {\large\bfseries Paper No. MTB-509: Special Theory of Relativity-I (2012)}\\[4pt]
  {\normalsize Time: 3 Hours \hfill Full Marks: 70}
\end{center}
\rule{\linewidth}{0.4pt}
\smallskip
\noindent\textit{Instructions: Answer \textbf{five} questions including Question~1 which is compulsory. Figures in right-hand margin indicate marks.}
\bigskip

\noindent\textbf{Question 1.}
\begin{parts}
  \item Show that Newton's second law of motion is invariant under special Galilean transformations. \pts{4}
  \item Show that the Lorentz transformation equations are superior to Galilean transformations. \pts{4}
  \item At what speed should a clock be moved so that it may appear to lose 1 minute in each hour? \pts{2}
  \item The length of a rocket ship on the ground is $100\,\text{m}$. When it is in flight, its length observed from the ground is $99\,\text{m}$. Calculate its speed. \pts{2}
  \item Show that if $V_0$ is the rest volume of a cube, then $V = V_0\sqrt{1 - v^2/c^2}$ is the volume viewed from a frame moving parallel to an edge of the cube. \pts{2}
  \item Show that the circle $x^2 + y^2 = a^2$ in frame $S$ appears as an ellipse in frame $S'$ moving with velocity $v$ along the $x$-axis. \pts{3}
  \item Explain Minkowski four-dimensional space-time. \pts{3}
  \item Define the 4-dimensional flat space-time metric. What is its use for 4-D vectors? \pts{2}
\end{parts}

\medskip\rule{\linewidth}{0.2pt}

\medskip\noindent\textbf{Question 2.}
\begin{parts}
  \item A Cartesian frame $OX'YZ$ rotates about the axis $OZ$ with uniform angular velocity $\omega$ relative to the inertial frame $OXYZ$. Show that $OX'YZ$ is non-inertial. \pts{6}
  \item Prove that the Lorentz transformation equations can be written in the matrix form:
  \[
  \begin{pmatrix} x' \\ ct' \end{pmatrix} = \begin{pmatrix} \cosh\theta & -\sinh\theta \\ -\sinh\theta & \cosh\theta \end{pmatrix} \begin{pmatrix} x \\ ct \end{pmatrix}
  \]
  Hence show that Lorentz transformations can be regarded as a rotation through an imaginary angle. \pts{6}
\end{parts}

\medskip\rule{\linewidth}{0.2pt}

\medskip\noindent\textbf{Question 3.}
\begin{parts}
  \item Prove that the set of all Lorentz transformations satisfies the group properties. \pts{6}
  \item Determine the component velocities of a particle with respect to two different frames $S$ and $S'$ using the Lorentz transformation (relativistic velocity addition). \pts{6}
\end{parts}

\medskip\rule{\linewidth}{0.2pt}

\medskip\noindent\textbf{Question 4.}
\begin{parts}
  \item By using Lorentz transformation equations, show that simultaneity is not absolute but relative. \pts{6}
  \item Calculate the length of a rod moving with velocity $0.8c$ in a direction inclined at $60°$ to its own length. The proper length is $100\,\text{cm}$. \pts{6}
\end{parts}

\medskip\rule{\linewidth}{0.2pt}

\medskip\noindent\textbf{Question 5.}
\begin{parts}
  \item Prove that the three-dimensional volume element $dx\,dy\,dz$ is not Lorentz invariant, but the four-dimensional volume element $dx\,dy\,dz\,dt$ is invariant under Lorentz transformation. \pts{6}
  \item If $u$ and $u'$ are velocities of a particle in two inertial systems $S$ and $S'$ (where $S'$ has velocity $v$ relative to $S$ in the $x$-direction), show that:
  \[
  u^2 = \frac{u'^2 + v^2 + 2u'v\cos\theta'- \left(\dfrac{u'v\sin\theta'}{c}\right)^2}{\left(1 + \dfrac{u'v\cos\theta'}{c^2}\right)^2}
  \]
  where $\theta'$ is the angle $u'$ makes with the $x$-axis. \pts{6}
\end{parts}

\medskip\rule{\linewidth}{0.2pt}

\medskip\noindent\textbf{Question 6.}
\begin{parts}
  \item Find the velocity at which the mass of a particle is double its rest mass. \pts{5}
  \item Explain the Michelson-Morley experiment. \pts{7}
\end{parts}

\medskip\rule{\linewidth}{0.2pt}

\medskip\noindent\textbf{Question 7.}
Explain the following terms in detail: (a) Light cone; (b) World points and world lines; (c) Space-like intervals; (d) Time-like intervals. \pts{12}

\vfill\rule{\linewidth}{0.4pt}
\begin{center}\small
  \href{https://bhu-exam.vercel.app/}{Click here to visit our website}\\[4pt]
  \href{https://me-aryan.vercel.app/}{Click here to visit our contributor}\\[4pt]
  \href{https://quashan-me.vercel.app/}{Click here to visit our contributor}
\end{center}
\end{document}
